% novices.tex                
% Note, that this a HyperLaTeX source and not plain LaTeX!
\documentclass[a4paper,11pt,twoside,titlepage,dvips]{scrartcl}
\usepackage{hyperlatex}
\usepackage{a4wide}
\usepackage{times}
\usepackage[latin1]{inputenc}
\usepackage[T1]{fontenc}
\usepackage{graphicx}
\usepackage{alltt}
\usepackage{moreverb}
\usepackage{fancyhdr}
\W\usepackage{rhxpanel}
\W\usepackage{sequential}


\T\DeclareGraphicsExtensions{.eps.gz,.eps}


% Hyperref should be among the last packages loaded
\usepackage{hyperref}

% Macros specific to this package
\newcommand{\Button}[1]{[\,#1\,]}
\newcommand{\Menu}[1]{\emph{#1}}
\newcommand{\Filename}[1]{\texttt{#1}}

% Some people still insist that Email should be written as two words
% for whatever reasons.  Maybe is is that they don't like the game of
% Teekesselchen and fear the ambiguity between email and enamel (which
% uses the same word in German).  

\newcommand{\Email}{email}

\newcommand{\EchelonUrl}{http://www.heise.de/tp/r4/artikel/6/6928/1.html}
\newcommand{\EinsteigerPDFURL}{http://wald.intevation.org/frs/download.php/134/gpg4win-fuer-einsteiger-2.0.0.RC1.pdf}
\newcommand{\DurchblickerPDFURL}{http://wald.intevation.org/frs/download.php/134/gpg4win-fuer-durchblicker-2.0.0.RC1.pdf}

\newcommand{\OtherBook}{\texorhtml{$\spadesuit$}{\xml{p}==\xmlent{gt}}}

\newcommand{\IncludeImage}[2][]{\texorhtml{%
\includegraphics[#1]{#2}%
}{%
\htmlimg{#2.png}%
}}

\newcommand{\HlxIcons}{}

%%\W\htmlpanelgerman
%%\W\extrasgerman
%%\W\dategerman
%%\W\captionsgerman
%%\htmlpanelfield{Inhalt}{hlxcontents}
\W\renewenvironment{thefootnotes}{\chapter{Footnotes}
\W  \begin{description}}{\end{description}}


% Get the version information from another file.  That file is created
% by the configure script.
\input{version.tex}



\newcommand{\manualversion}{\manualversionNovices}
\newcommand{\manualdate}{\manualdateNovices}

\T\fancyhead{} % clear all fields
\T\fancyhead[LO,RE]{Gpg4win for Novices \manualversion\ \manualinprogressEn}
\T\fancyhead[RO,LE]{\thepage}
\T\fancyfoot[C]{\includegraphics[width=1cm]{gpg4win-logo}}


% Title stuff 
\htmltitle{Novices}
%\htmladdress{Gpg4win Project, \today}
\title{
\IncludeImage[width=8cm]{gpg4win-logo}
\\
Gpg4win for Novices}

\author{\htmlonly{\xml{p}\small
%%Here \xlink{Download PDF Version 2.0.1 (1.5
%MByte)}{http://wald.intevation.org/frs/download.php/203/gpg4win-fuer-einsteiger-2.0.1.pdf} und
\xlink{List of PDF Download
Versions}{http://wald.intevation.org/frs/?group_id=11}\xml{br}
To \xlink{Gpg4win for Advanced Users}{advanced.html}\xml{br}
To \xlink{Gpg4win Homepage}{http://www.gpg4win.org/}\xml{p}
}%
A publication of the Gpg4win project\\
  \small based on original documents by  
\T\\
  \small Manfred J. Heinze, Karl Bihlmeier, Isabel Kramer,
\T\\
  \small Dr. Francis Wray, Ute Bahn,
\T\\ \
  \small Werner Koch.
 \\
  \small Translated from the German original by 
\T\\
  \small Brigitte Hamilton}
\date{Version \manualversion\ as of \manualdate\ \manualinprogressEn}


\begin{document}
\thispagestyle{empty}
\pagestyle{fancy}
\T\parindent0cm
\T\parskip\medskipamount


\maketitle


\section*{Imprint Gpg4win}

\noindent
Copyright \copyright{} 2002 Bundesministerium f�r Wirtschaft und
Technologie (German Federal Ministry of Economics \& Technology)\\
Copyright \copyright{} 2005 g10 Code GmbH\\
Copyright \copyright{} 2006 Brigitte Hamilton\\
Permission is hereby granted to copy, distribute and/or modify this
document under the terms of the GNU Free Documentation License, 
Version 1.2 (or newer) published by the Free Software Foundation; 
with the Invariant Sections being "Impressum", no Front-Cover Texts, and
no Back-Cover Texts.  A copy of the license is included in the section
entitled "GNU Free Documentation License".

%%\htmlonly{Die aktuelle PDF Version dieses Dokuments finden sie unter
%%\xlink{\EinsteigerPDFURL}{\EinsteigerPDFURL}.}

Special Thanks to:\\
Bernhard Reiter for minor edits and coordination (2006, 2007).\\
Michael W. Lewis for proofreading (2007).\\
Andre Hosch for screenshots (2007)

\vfill

\begin{center}
  \T\fbox{\parbox{.7\textwidth}{%
      The statements on the \textbf{next page} are not anymore
      correct.  Due to a wrong application of the terms of
      the GFDL, it is legally not possible to fix them.  Please add
      new copyright notices only here.  \T}}
\end{center}

\clearpage
%% Orginal page 4
%% This page 4 is an invariant section. Translation not allowed.
\section*{Impressum}
Diese Seite darf nicht ver�ndert werden.\\

Autor: Manfred J. Heinze, TextLab text+media\\
Beratung: Lutz Zolondz, G-N-U GmbH\\
Illustrationen: Karl Bihlmeier, Bihlmeier \& Kramer GbR\\
Layout: Isabel Kramer, Bihlmeier \& Kramer GbR\\
Fachtext: Dr. Francis Wray, e-mediate Ltd.\\
Redaktion: Ute Bahn, TextLab text+media\\
1. Auflage, M�rz 2002\\


Copyright \copyright{} Bundesministerium f�r Wirtschaft und Technologie

Dieses Buch unterliegt der "`GNU Free Documentation License"'.
Originaltext der Lizenz: http://www.\texonly{\linebreak{}}gnu.org/copyleft/fdl.html.
Deutsche �bersetzung http://nautix.sourceforge.net""/docs/fdl.de.html
sowie auf der beiliegenden CD-ROM.  Es wird die Erlaubnis gegeben,
dieses Dokument zu kopieren, zu verteilen und/oder zu ver�ndern unter
den Bedingungen der GNU Free Documentation License, Version 1.1 oder
einer sp�teren, von der Free Software Foundation ver�ffentlichten
Version.  Diese Seite ("`Impressum"') darf nicht ver�ndert werden und
muss in allen Kopien und Bearbeitungen erhalten bleiben
("`unver�nderlicher Abschnitt"' im Sinne der GNU Free Documentation
License).  Wenn dieses Dokument von Dritten kopiert, verteilt und/oder
ver�ndert wird, darf in keiner Form der Eindruck eines Zusammenhanges
mit dem Bundesministerium f�r Wirtschaft und Technologie erweckt
werden.  Wie das OpenSource-Kryptografieprogramm GnuPP selbst wurden
diese Texte nicht f�r Mathematiker, Geheimdienstler und Kryptografen
geschrieben, sondern f�r jedermann.



\clearpage %% End of original page 4.

\tableofcontents

%%\clearpage
%% Orginal page  6
%% We don't use these foreword anymore because Mr. M�ller is not
%% anymore minister of economic and technology.  We might want to ask
%% for a new foreword by the current head of that ministr
%%

\clearpage
%% Orginal page 7
\section{About this manual}

The Gpg4win manual and exercise module consists of three parts:

\begin{itemize}

\item \textbf{an introductory guide called "Gpg4win for Novices",} 
which you are reading right now,

\item \textbf{the manual "Gpg4win for Advanced Users"} in PDF-Format,
which can be found on your hard drive after you have installed Gpg4win,

\item \textbf{the exercise robot Adele,} which allows you to practice
the \Email{} en- and decryption process(an internet connection is
required). 
\end{itemize}


\textbf{"Gpg4win for Novices"} is a quick guide to the installation and
everyday use of the Gpg4win software. You will need about half an hour to work
through this manual, depending on your knowledge of computers and Windows. 

\textbf{"Gpg4win for Advanced Users"} provides in-depth information
about the basic principles and mechanisms usbed by Gpg4win, 
as well as its less commonly used capabilities.

Both manuals are available in PDF format; you can print your own
manual if you did not receive a printed copy. 

Each manual can be read independently. We suggest, however, that you 
read both manuals in order to get a better understanding of the software.

\OtherBook{} This symbol references a link to the other manual.


\textbf{The practice robot Adele} is available on the Internet. Adele
receives, sends and decrypts encrypted \Email{}s. You can use it to 
practice a complete cryptographic exchange, as often as you need 
to be fully familiar with the software.

Adele was developed as part of the older GnuPP project, where it is
still used.  "Gpg4win for Novices" also uses this very reliable
practice robot and hereby wishes to express its gratitude to the owners
of gnupp.de for operating Adele. 

\clearpage
%% Orginal page 8
\section{What is Gpg4win?}

% what-is-gpg4win.tex  - Introduction used by both manuals


\textbf{
Gpg4win (GNU Privacy Guard for Windows) is an email encryption software.
It is the result of a project initiated by the Federal Office for
Information Security, and includes the following components: 
}

\textbf{GnuPG}: its key component, the encryption software \newline
\textbf{GPA}: GNU Privacy Assistant,
a key manager \newline
\textbf{WinPT}: Key Manager, which also supports encryption via your
Clipboard \newline
\textbf{GPGol}: a plugin for Microsoft Outlook which integrates the
operation of GnuPG \newline
\textbf{GPGee}: a plugin for Windows Explorer which allows encryption of
data by right-clicking on your mouse \newline
\textbf{Claws Mail}: a complete email program with integrated GnuPG
operation \newline

The encryption program GnuPG (GNU Privacy Guard) provides you with a
secure, simple and free method of email encryption. It can be used
privately or commercially without restrictions.  The encryption
technology used by GnuPG is extremely secure and cannot be broken using
current technology.

GnuPG is a free software\footnote{sometimes incorrectly identified as
Open Source Software}. This means that anyone can use the software for
private or commercial purposes, as well as analyze or change the source
codes (ie. the actual programming commands), and distribute the
same.\footnote{However, keep in mind that a fair amount of technical
knowledge is required to change the program, otherwise the program's
security may be compromised.}

The transparency of the source code forms an essential part of a
security software, as it is the only way to verify the trustworthiness
of the program.


GnuPG is based on the international standard OpenPGP (RFC 2440), is
fully compatible with PGP and uses the same infrastructure (key server
etc.). Since GnuPG version 2 the cryptography standard
S/MIME (CMS/RFC 3852 and/or X.509) is also supported.

PGP ("`Pretty Good Privacy"') is not free software; many years ago it
was available on a temporary basis under similar conditions as GnuPG,
but is no longer considered state-of-the-art.

Additional information regarding GnuPG and other projects undertaken by
the german government in the area of Internet security can be found on
the website of the Federal Office for Information Security \W
\xlink{www.bsi-fuer-buerger.de}{http://www.bsi-fuer-buerger.de}
\T\href{http://www.bsi-fuer-buerger.de}{www.bsi-fuer-buerger.de}





\clearpage
%% Orginal page 9
\section{Installing Gpg4win}

If you already have a GnuPG-based application such as GnuPP, GnuPT,
WinPT or GnuPG Basics installed on your system, we recommend that you
read the Appendix \ref{ch:migration} to find out how to migrate existing
keys. 

Installing Gpg4win from a CD-ROM:

Insert the CD-ROM in the CD-ROM drive of your computer and log in as
Administrator. On the screen click on the CD-ROM icon titled 'Gpg4win'.
Once the CD-ROM icon opens, click on the installation icon titled
'Gpg4win'.

Installing Gpg4win from the Internet:

If you have downloaded Gpg4win from the Internet, click on the new file
(it should be named \texttt{gpg4win-\PackageVersion{}.exe} or a newer
version).  Please ensure that you have downloaded the file from a
trustworthy site.

The following installation steps apply to all situations:

You will be asked if you want to install the program; click on
\Button{Yes}. 

The following screen will appear:

% screenshot:  Welcome Seite Installer
\begin{center}
\IncludeImage{sc-en-inst-welcome}
\end{center}

If you have other programs running on your computer, close them now and
click on \Button{Next}.

\clearpage
%% Orginal page 10

The licencing page contains information regarding the licencing of this
software.

If your sole intention is to install and use the software, you do not
need to read this information. 

However, if you plan to distribute and/or alter the software, you must
familiarize yourself with the conditions contained in the licencing
agreement.


% Screenshot Lizenzseite des Installers
\begin{center}
\IncludeImage{sc-en-inst-license}
\end{center}


Click on \Button{Next}.


\clearpage
%% New page (not in original document)

On the components page you can select which features of GnuPG For
Windows you want to install. 
 
To assist you with your selection, a short description window appears
when you move the mouse over a selected item. At this point, you may
want to check the available free space on your hard drive as well.

% sreenshot Auswahl zu installierender Komponenten
\begin{center}
\IncludeImage{sc-en-inst-components}
\end{center}

The recommended minimum installation consists of GnuPG, GPA, WinPT and
the manuals. If required, the remaining programs can be installed later.

Click on \Button{Next}.


\clearpage
%% Original page 11

On this page you can choose the folder in which to install Gpg4win on
your computer. If you do not enter a folder name, 
you should accept the default folder name, which is: \\
\Filename{C:$\backslash$Programme$\backslash$GNU$\backslash$GnuPG}\\%%FIXME

% screenshot: Auswahl des Installationsverzeichnis.
\begin{center}
\IncludeImage{sc-en-inst-directory}
\end{center}

Click on \Button{Next}.

\clearpage

This page allows you to set start links for the program. The default
setting adds Gpg4win to the start menu only.  You can change
these settings within Windows at anytime.

% screenshot: Auswahl des Links
\begin{center}
\IncludeImage{sc-en-inst-options}
\end{center}

Click on  \Button{Next}.

\clearpage
%% Original page 12

If you chose to add the program to the Start Menu (as per previous
page), this page allows you to choose a Start Menu subfolder for the
program.


% screenshot:  Startmenu ausw�hlen
\begin{center}
\IncludeImage{sc-en-inst-startmenu}
\end{center}

For a standard installation, select the default setting and click on
\Button{Install}.

\clearpage


During the installation a progress bar displays the file being installed.
You can click on \Button{Show details} to see a log of the installation.

% Screenshot: Ready page Installer
\begin{center}
\IncludeImage{sc-en-inst-ready}
\end{center}

Once the installation is complete, click on \Button{Next}.

\clearpage
%% Original page 13

The following page shows the last step of the installation process: 

% Screenshot: Finish page Installer
\begin{center}
\IncludeImage{sc-en-inst-finished} 
\end{center}

Click on  \Button{Finish}.

\clearpage

You may need to restart Windows for the settings to take effect.  In
that case, instead of showing the installation completion page, the
following page appears: 

% Screenshot: Finish page Installer with reboot
\begin{center}
\IncludeImage{sc-inst-finished2}
\end{center}

At this point you can choose to restart Windows automatically or restart
later manually.

Click on \Button{Finish}.


% FIXME:  Wir m�ssen erkl�ren wie man Word als Standard Editor in
% Outlook ausschaltet.

\clearpage
%% Original page 14

\textbf{And that's it!}


You have now installed Gpg4win and are ready to use the program.

Prior to starting Gpg4win, we recommend reading Chapter 3 and 4 of the
manual "Gpg4win for Advanced Users"
(PDF-File).  These chapters highlight the ingenious theory behind
Gpg4win's ability to encrypt \Email{}s in a safe and user-friendly
manner.

Of course you do not need to know all the technical details of Gpg4win to be
able to use it. However, since you will use Gpg4win to handle your 
most sensitive correspondence and it probably iis a good idea to understand
the theory behind it.

$\ldots$

The following section provides you with tips for creating a
secure yet easy-to-remember passphrase.

\textbf{
\OtherBook{}  At this point, please read Chapter 3 and 4 in the manual
"\xlink{Gpg4win for Advanced Users}{durchblicker.html}" before reading
on.}

\clearpage
%% Original page 15
\section{Creating a key pair}

After reading the information explaining Gpg4win's security features
and the creating of a good passphrase to protect your
private key, this section shows you how to create a key pair.

\textbf{
The processes of creating a key, encryption and decryption are very
important - so important that it should be possible to practice $\ldots$
}

\textbf{And this is what you actually can do:}
Run through the whole processes as often as you want.
Those "dry-runs" help you gain confidence in
using the program, so that some of the more intricate steps in creating
keys will not pose any problems later on. You can use Adele to do this.
It is a test server which was developed as part of the GnuPP project.
Adele allows you to try and test several practice keys before creating
your own set.  

\clearpage
%% Original page 16
\textbf{Let's get started!}

Start the GPA program from your Windows Start menu: 

% screenshot Startmenu with GPA highlighted
\begin{center}
\IncludeImage[width=0.4\textwidth]{sc-en-gpa-rungpa}
\end{center}

The following window appears: 

% screenshot: GPA erste Schl�sselerzeugung
\begin{center}
\IncludeImage[width=0.9\textwidth]{sc-en-gpa-nokey}
\end{center}

Click on  \Button{Generate key now}.


\clearpage
%% Original page 17

\textbf{For practice purposes} you can enter any name in the window for
now, e.g. "Larry Smith".

% screenshot: Key generation wizard first page
\begin{center}
\IncludeImage[width=0.6\textwidth]{sc-en-gpa-gen-name}
\end{center}

Or you can generate your 'real' key in which case you should enter your
real name.

When you are done, click on  \Button{Forward}.

\clearpage
%% Original page 18
Now enter your \Email{} address.

Again, if this is a test run, you can use an imaginary \Email{} address
such as 
"\verb-larrys@gpg4win.de-"

% screenshot: Wizard email address
\begin{center}
\IncludeImage[width=0.6\textwidth]{sc-en-gpa-gen-email}
\end{center}

Or you can enter your real \Email{} address, and click on
\Button{Forward}.

%note by translator Brigitte: when I ran Gpg4win, the option outlined in
%the text below was not available!  FIXME

This option allows you to enter notes regarding your key.  Usually this
field is empty; however, if you are creating a test key, you should make
a note, such as "test". The notes are a part of your User-ID and will
be made public along with your name and \Email{} address.  Now click on
\Button{Forward}.


\clearpage
%% Original page 19

Entering a passphrase is one of the most important steps in the program,
as the program is only as good (and as secure) as your passphrase! 

Chapter 4 ("The Passphrase") from the manual "Gpg4win for Advanced
Users" already provided you with suggestions on creating a secure
passphrase.  So, at this point you should have come up with your own
passphrase - one that is private, easy to remember, and difficult to
crack.  

Enter your passphrase now. 

% screenshot: Wizard, Enter passphrase
\begin{center}
\IncludeImage[width=0.6\textwidth]{sc-en-gpa-gen-passwd}
\end{center}

If the passphrase you entered is not very secure, 
a warning will be displayed,
allowing you to enter a more secure passphrase. 

\textbf Again, you can enter a test passphrase for practice purposes, or
your real passphrase.

Once you have entered your passphrase twice, click on \Button{Next}.

\clearpage
%% Original page 21

This starts the creation of your key pair, which can take a few minutes.
In the meantime, you can continue to work with other programs on your
computer which will slightly increase the quality of the key being
generated.

%put a note in here somewhere that 'key pair' and 'key' are used
%interchangeably.  FIXME

Once the key is generated, the following window appears:

% screenshot: Wizard, backup
\begin{center}
\IncludeImage[width=0.6\textwidth]{sc-en-gpa-gen-backup}
\end{center}

This window asks you to create a back-up copy of your key. Please do
this now, even if this is just a test run.

If the default file name is acceptable, click on  \Button{OK}. If you
would like to store the back-up in a different location, please select a
different file name now. 

\textbf{Important:} Once the the back-up file is saved onto your hard
drive, you should save this file onto a removable storage device
(e.g. USB stick, diskette or CD-ROM), and subsequently delete the
original file.  Be sure to keep the storage device in a secure place. 

Please note that you can create back-up copies at any time by selecting
the following option from the main menu: 
\Menu{Key$\rightarrow$Backup}.


\clearpage
%% Original page 23
\textbf{
This concludes the installation of Gpg4win and the generation of your
key pair.  You are now the owner of a unique and secure digital key. 
}

You should now see the GPA's main window. The key pair you just created
will be shown in the middle of the window, beside the key pair symbol.  

Clicking on the key pair symbol will provide details about your key
pair, which are discussed in more detail below.

% screenshot: keymanger mit einem Schl�ssel.
\begin{center}
\IncludeImage[width=0.9\textwidth]{sc-en-gpa-first-key}
\end{center}

What is the significance of the notes regarding your key? Your key has
no built-in expiry date and is, therefore, valid indefinitely.  You can,
however, change the key's validity period - more about this later.

A key consisting of 1024 bits is considered very secure, without
creating undue strain on your computer's capacity.


\textbf{ \OtherBook{} More information on this topic can be found in
Chapter 5 "More about keys" in the manual "\xlink{Gpg4win for
Advanced Users}{durchblicker.html}".}

\clearpage
%% Original page 24

\section{Publishing your key per  \Email{}}

One of the more practical aspects of Gpg4win is its ability to use a
"non-secret" public key for the en- and decryption of data.  As long
as your key and its corresponding passphrase are secure, you have gone a
long way towards keeping your information confidential. 

Everyone can and should have your public key, just as you can and should
have the public keys of the people you correspond with. 

Because:

\textbf{
In order to exchange secure \Email{}s, each party must have access to
the public key of the other party. 
}

In order to send someone encrypted \Email{}s, you must have their
public key in order to encrypt the emails to sent to them. 

Similarly, someone wanting to send you encrypted \Email{}s must 
have your public key in order to encrypt the email being sent to you.  

This is the reason that your public key should be made widely
available. Depending on the number of correspondence partners,
can do this in two ways:

\begin{itemize}
\item \textbf{directly sending an \Email{}} to selected recipients
\item \textbf{or publishing the key on a key server} --- making it
available to anyone
\end{itemize}

The first way to distribute your public key is to send it by
\Email{} to one or more selected recipients.
Alternatively you can make your \Email{} address available to anyone
in the Internet. The second option is somewhat risky, as it can
result in considerable SPAM activity on your email account. Therefore it
is a good idea to only use an address with a good SPAM filter. 


\clearpage
%% Original page 25
You can use Adele to practice the following steps:

Adele is a very good \Email{} robot for practicing secure correspondence.
Because most people prefer to correspond with a real person rather
than with a piece of software (which is what Adele is, after all), we
developed the following scenario:

% Cartoon:  Adele mit Buch ind er Hand vor Rechner ``you have mail"'
\begin{center}
\IncludeImage{adele01}
\end{center}

You first send Adele your public key. Once Adele has received your key,
she uses it to encrypt an  \Email{} which she sends to you along with
her own public key.  You can now decrypt Adele's email using your own
private key, and you can also respond to Adele by encrypting your email
with her public key.

You can now export your public key, copy it into an \Email{} und send it
to Adele.


\clearpage
%% Original page 26

Here is one possible way of doing just that, a method which works even
if your \Email{} service does not allow attachments.
This procedure also gives you a first in-depth look at your
key and its components.

How it works:\\
Select the key you want to export by clicking on the corresponding key
on your list, and then clicking on the \Button{Export} icon of the main
GPA menu. Choose a file to export your key to, e.g. \\
\Filename{my-key.asc}.  A popup window will let you know whether the
operation was successful.  Then click on \Button{OK}.


\clearpage
%% Original page 27

You can access the file through Windows Explorer; make sure you choose
the same folder you chose when exporting the key.  You can open
the file with a text editor (ex. WordPad), which will show your public
key as a series of blocks containing text and numbers. 

% screenshot: Editor mit ascii armored key
\begin{center}
\IncludeImage[width=0.9\textwidth]{sc-misc-mein-key-asc}
\end{center}


\clearpage
%% Original page 28

Highlight the whole key portion from 

\begin{verbatim}
-----BEGIN PGP PUBLIC KEY BLOCK-----
\end{verbatim}
to
\begin{verbatim}
-----END PGP PUBLIC KEY BLOCK-----
\end{verbatim}

and copy it using the copy function on your toolbar or a keyboard
shortcut such as Ctrl-C. This saves your key on the clipboard until you
are ready to paste it, as described below. 

Start your email program. Open a new \Email{} message and paste
your public key (Windows users may use a shortcut key such as Ctrl-V).
Before doing this, you should configure your email program to
send messages in text-only format rather than HTML.

Put \verb-adele@gnupp.de- into the address line of the \Email{}, and
\textbf{my public key} on the subject line.

Your \Email{} should look like this:

% screenshot:  Eines composer Windows.
\begin{center}
\IncludeImage[width=0.9\textwidth]{sc-ol-send-test-key}
\end{center}

Now send this \Email{} to Adele. Make sure that you include your proper
email address as sender rather than your practice address, otherwise you
will never get a response from Adele!

\clearpage
%% Original page 29

This process works exactly the same way if you send your key to a real
\Email{} address. You can add other text, just like in any other email.
Obviously, this is not required for emails addressed to Adele, as the
robot's only purpose is to help you with the technical aspects of this
process.

\textbf{Summary:} You have now sent your public key by \Email{} to
someone else (e.g. Adele). 

\textbf{ \OtherBook{} The "Copy \& Paste" method shown in this
example is easily understood by beginners.
Chapter 7 of the manual "\xlink{Gpg4win for Advanced
Users}{durchblicker.html}" describes how to send your key as a
file attachment, which is a more commonly used method.}  

\clearpage
%% Original page 30
\section{Sending your key to a keyserver}

You can use this option even if you are exchanging encrypted \Email{}s
with only a few people. By making your key available on a keyserver, it
is always accessible to you and/or others.  

\textsc{PLEASE NOTE: Although there are no indications that spammers are
currently collecting addresses from keyservers, it is technically
possible to do so. We do not, therefore, recommend publishing your key to
a key server if you do not have an effective spam filter.}

Choose your key by clicking on it. Then click on \Menu{Server
$\rightarrow$Send key} which is found under \Menu{Server$\rightarrow
$Server}.  A popup window will alert you that your key is about to be
sent to a public key server. 

% screenshot: GPA export key to keyserver
\begin{center}
\IncludeImage[width=0.4\textwidth]{sc-en-gpa-ks-export-p}
\end{center}

This window also displays a default key server.  If you click on
\Button{Yes}, your key will be automatically sent to the server, from
where it is sent to a network of worldwide key servers. Now anyone can
download your key, and use it to send you a secure \Email{}.

If this is a practice run, do NOT send your practice key to the key
server because once sent, it cannot be removed.

\textbf{Summary:} You now know how to distribute your key via a key
server on the Internet. 

\textbf{ \OtherBook{} Chapter 6 of the manual"\xlink{Gpg4win for
Advanced Users}{durchblicker.html}" explains how to find someone else's
key on a keyserver.}

\clearpage
%% Original page 31
\section{Decrypting an \Email{}}

Adele has received your public key, and uses it to encrypt an \Email{}
which she sends back to you. 

% cartoon: Adele typing and sending a mail
\begin{center}
\IncludeImage{adele02}
\end{center}


\clearpage
%% Orginal page 32
The email will look like this:

\begin{verbatim}
From: Adele (The friendly email-Robot) <adele@gnupp.de>
Subject: Re: my publis key
To: larrys@gpg4win.de
Date: Thu, 12 Jan 2006 09:17:28 +0100

-----BEGIN PGP MESSAGE-----
Version: GnuPG v1.4.1 (GNU/Linux)

hQEOA9FS8I3hSvdPEAP/W6W6f4MBwqTdzd9O/7FOTDHh//bQ+GUWoT0k9Y0i96UZ
QO1VhQSia6a8DZrFQj7SlJWmB1MM7RNhkmhfZsD5Bn9ICmwwOt2xJDBkCQ34gu5N
NxQ92WXZjHCaI0dSlynNziNbK8Ik26YPBYkQjLUDhHN4CRZ7q67eVEd/B9DI04wD
....
ujbjyj09L/9NvoBniWrgqVUayKr1Ls8OIZkyiex6mKypPGADJFAzvTwjubj5S6zJ
A+QvSXUB9Hj8Ft2Nt3j0B/gWn5no3Er2/15UcBn/UPSxW9or0w9seDxCuSXvpakX
bcneOm/pcJNEHcApXWXpoNOxRZ1MksM300w+79M6p2w=
=VCHb
-----END PGP MESSAGE-----
\end{verbatim}

(Please note that the encryption sequence has been considerably 
shortened to save space in this manual.)
\clearpage
%% Original page 33

\textbf{Using WinPT to decrypt this \Email{}.}

WinPT (Windows Privacy Tray key management tool) serves as PnuPG's
"front end". This program is used to encrypt and decrypt \Email{}s; it
also creates and verifies digital signatures.  One of WinPT's advantages
is that it works with all email programs. 

Most email programs (e.g. MS Outlook for Windows) also have special
plug-ins which can encrypt and decrypt \Email{} directly within the program.


\textbf{ \OtherBook{} Chapter 8 of the manual "\xlink{Gpg4win for
Advanced Users}{durchblicker.html}" provides more information on this
particular funtion.} 

WinPT performs the en/decryption function using your computer's memory.
This means that any text to be en/decrypted must be copied onto the
computer's clipboard.

To do this, highlight the complete text in Adele's \Email{} and copy it
onto the Clipboard using the copy function or the shortcut Ctrl-C. 


\clearpage
%% Original page 34
Starting WinPT from Windows Start Menu:

% screenshot: startmenu Auswahl winpt
\begin{center}
\IncludeImage[width=0.7\textwidth]{sc-winpt-startmenu}
\end{center}

% the english version of the program does not have the "hinweis"below,
% so I have not translated it. FIXME

% W�hrend WinPT startet, erscheint kurz ein Hinweis darauf, dass das
% Programm die bereits vorhandenen Schl�ssel einl�dt.

Once you open the program, the WinPT icon (a key) appears in the
Windows taskbar, which is located in the lower right hand corner of the screen.

% screenshot: taskbar rechts mit WinPT icon
\begin{center}
\IncludeImage{sc-winpt-trayicon}
\end{center}


\clearpage
%% Original page 35
Right-click on the icon to open the menu box, and click on
\Menu{Clipboard$\rightarrow$Decrypt/Verify}.

The program will then ask you for your secret passphrase and proceed to
decrypt Adele's \Email{}.

% screenshot: WinPT decryption dialog
\begin{center}
\IncludeImage[width=0.7\textwidth]{sc-winpt-clip-decrypt}
\end{center}

A popup window window will notify you when the decryption process is
complete. 

\clearpage
%% Original page 36

The decrypted text is now located on the clipboard, much in the same way
as was done during the encryption process. You can copy the text (using
shortcut key Ctrl-V) into your text editor or email program. 

Adele's decrypted email will look something like this\footnote{Depending
on Adele's software version, it might look slightly different.}:

\begin{verbatim}
Hello Larry Smith,

here is an encrypted response to your email. 

I received your public key with the ID 57251332CD8687F6 and a
description `<larrys@gpg4win.de>'.

I have enclosed the public key of adele@gnupp.de.

Sincerely,
adele@gnupp.de
\end{verbatim}

The text sequence following this message is Adele's public key. 

The next section shows you how to import someone's public key and
attach it to your key ring, which allows you to use this key anytime
to encrypt messages to that person (in this case, Adele) or to verify
a digital signature.


\textbf{Summary:}

\begin{enumerate}
\item You know how to decrypt an encrypted \Email{} using your private
key.

\item You know how to send an encrypted response using the public key of
the sender. 
\end{enumerate}


\clearpage
%% Original page 37
\section{Attaching a key to your key ring}

You can store your friends' public keys on your "key ring" so they do
not have to send you their public key every time they write to you.  

\textbf{First Option:}

To import a public key (ie. attach it to your key ring), you can save
the key as a text block, much like you did with your own key. 


To do this:

Highlight the public key portion of the email you received, from

\begin{verbatim}
-----BEGIN PGP PUBLIC KEY BLOCK-----
\end{verbatim}
to
\begin{verbatim}
-----END PGP PUBLIC KEY BLOCK-----
\end{verbatim}

and use Copy \& Paste to insert it into your text editor.  We recommend
using a file and folder name that is easily found later; e.g.
\Filename{adeles-key.asc} in folder 
\Filename{My Documents}.



\clearpage
%% Original page 38
\textbf{2nd Option:}

The key is sent to you as a file attachment to the \Email{}. No matter
which mail program you use, you can always save attachments onto your
hard drive.  Do this now (again using names you will easily recognize
and find later on, e.g. \Filename{My Documents}).

It does not matter whether you save the key as text or directly as an
\Email{}attachment, as both methods import the key into your GnuPG-"Key
Ring".

% screenshot:  Startmenu, Auswahl GPA

This is how it works:

Start the GNU Privacy Assistant (GPA) from Windows (this is necessary
only if you shut it down after the previous practice session).


\clearpage
%% Original page 39
Click on Import, then select and load the key file.

% screenshot GPA Schl�sselverwaltung mit zwei Schl�sseln
\begin{center}
\IncludeImage[width=0.9\textwidth]{sc-en-gpa-two-keys}
\end{center}

You have now imported someone else's (in this case Adele's) public key
and attached it to your key ring. Now you can use this key to send
encrypted messages to the owner of that public key, as well as to verify
his or her signature.


\textbf{Before continuing, it is important to address the following
concern:}\\
It is possible that the \Email{} was sent by someone else using Adele's
name, therefore how do you know that the public key sent to you is
really Adele's key? 

\textbf{ \OtherBook{} Chapter 9 ("Key Verification") in the 
"\xlink{Gpg4win for Advanced Users}{durchblicker.html}" manual deals
with this important question.  You may want to read that section now 
before continuing with this manual. }

\clearpage
%% Original page 40
Chapter 9 of the manual "Gpg4win for Advanced Users" shows you how to
validate a key as well as how to sign a message (i.e. attach a signature)
using your private key.

Chapter 10 of the advanced manual also discusses ways to attach a
signature to \Email{} messages. This is the equivalent of attaching an
electronic seal to your message, allowing the recipient to verify
whether the \Email{} has been altered during transmission (and that it
came from you).

The signature verification process is simple. For this, you
need the sender's public key on your Gpg4win-"key ring" (see
Chapter 8 of "Gpg4win for Advanced Users" for more information).

% cartoon:  M�ller mit Schl�ssel
\begin{center}
\IncludeImage[width=0.4\textwidth]{man-with-signed-key}
\end{center}


\clearpage
%% Original page 41
You can tell whether an \Email{} has been electronically signed if the
text of the email is framed with the sender's signature (like a border).
It will look something like this: 

\begin{verbatim}
-----BEGIN PGP SIGNED MESSAGE-----
Hash: SHA1
\end{verbatim} 

and ends \Email{}-message with 

\begin{verbatim}
-----BEGIN PGP SIGNATURE-----
Version: GnuPG v1.4.2 (MingW32)

iEYEARECAAYFAjxeqy0ACgkQcwePex+3Ivs79wCfW8u
ytRsEXgzCrfPnjGrDDtb7QZIAn17B8l8gFQ3WIUUDCMfA5cQajHcm
=O6lY
-----END PGP SIGNATURE-----
\end{verbatim}

Highlight the text starting from \textsl{BEGIN PGP SIGNED MESSAGE} to
\textsl{END PGP SIGNATURE} and copy it (using Ctrl-C) to your
clipboard. 

Now continue to decrypt the \Email{} as shown in Chapter 7 of this
manual.

Right-click on the WinPT icon on your Windows taskbar and select
\newline
\Menu{Clipboard$\rightarrow$ Decrypt/Verify}.

You should see the following window:

% screenshot: WinPT good signature.
\begin{center}
\IncludeImage[width=0.7\textwidth]{sc-winpt-good-sig}
\end{center}

If the status line on the window displays the following message
\emph{Invalid Signature}%someone check if this is the right text used in
%the program, i did not encounter it so I am not sure of the tranlsation,
%although i think this should be right FIXME

it means that the message has been altered after being sent.  This does
not always mean that a third party has altered the message, as it could
also have been altered by a technical error during transmission
through the Internet. 

\textbf{ \OtherBook{} Before continuing, you may want to read Chapter 10
of the manual "\xlink{Gpg4win for Advanced Users}{durchblicker.html}"
which contains additional information on how to deal with
invalid signatures.}

\clearpage
%% Original page 42
\section{Encrypting \Email{}s}

Encrypt an \Email{} and send it to Adele (or another person).

Open your \Email{} program and write a message (if you are sending the
message to Adele, any text will do as Adele can not actually read $\ldots$)

Highlight the text and copy it onto your Clipboard (Windows) using the
copy function or the Ctrl-C shortcut.

% screenshot: composer mit mail an Adele, body markiert
\begin{center}
\IncludeImage[width=0.9\textwidth]{sc-ol-send-enc-msg1}
\end{center}

Right click on the WinPT icon on your Windows Taskbar and select
\Menu{Clipboard$\rightarrow$Encrypt}.


\clearpage
%% Original page 45
You should now see a window containing the keys on your key ring. Using
the examples featured in this guide, you would see Adele's key (the one
she sent to you), as well as your own key, which you created in Chapter
2 (FIXME: Use dynamic reference).

% screenshot: Winpt encryption key selection
\begin{center}
\IncludeImage[width=0.7\textwidth]{sc-winpt-sel-enc-key}
\end{center}

Click on Adele's key to encrypt the message you just composed, keeping in
mind the following principle:

\textbf{In order to send someone encrypted \Email{}s, you need to have
that person's public key in order to encrypt the message properly.}

Clicking on \Button{OK} starts the encryption process, which will be
confirmed by the program if successful.

\clearpage
%% Original page 46
The encrypted message can be found on your computer's clipboard, from
which you can copy it into the \Email{} window.  You can erase the
unencrypted text or just copy and paste the contents of the clipboard.

It will look something like this: 

% screenshot: composer Windows mit enrcypted text 
\begin{center}
\IncludeImage[width=0.9\textwidth]{sc-ol-send-enc-msg2}
\end{center}

Now send your \Email{} to Adele.  Remember to use your proper email
address (rather than the previous practice email) otherwise you will not
receive a response.   $\ldots$

\textbf{
Congratulations! You have just encrypted your first \Email{}!
}

\clearpage
%% Original page 47
\section{How to archive/store encrypted \Email{}s}

It is probably not wise to store the readable (clear)text of your
encrypted \Email{}s on your computer (after all, they were encrypted 
for a reason). Therefore it is sensible to store encrypted
copies of the \Email{}s.

You may already have guessed resulting problem: To decrypt the
archived \Email{}s you need the private key of the recipient, which will
(should) never be available to you. $\ldots$

The solution: You encrypt the messages to yourself as well. 

You can encrypt a message designated for a recipient (e.g. Adele) using
their own as well as your own public key, which allows you to decrypt the
text later using your own private key.

Since Gpg4win cannot tell the difference between the different keys used
to encrypt a message (ie. you could have more than one key), you need to
set up the program accordingly. 

\clearpage
%% Original page 48
To use this option, right click on the WinPT icon select
\Menu{Preferences $\rightarrow$ GPG}.

% screenshot: Winpt configuration dialog
\begin{center}
\IncludeImage[width=0.7\textwidth]{sc-winpt-enctoself}
\end{center}

This opens the GnuPG Preferences Window. You can add your key in the
"Encrypt to this key" field using your \Email{}address.

\Email{} programs which directly support GnuPG will also feature this
option.

\clearpage
%% Original page 49
\textbf{Summary:}

\begin{enumerate}
\item You responded to the recipient by encrypting an \Email{} with
his/her public key.

\item You also set up WinPT to encrypt archived copies of your \Email{}s
with your own public key.
\end{enumerate}

\textbf{And that's it - welcome to the world of free and secure \Email{}
encryption!}

\textbf{
\OtherBook{} We recommend you read Chapter 10 to 12 of the manual "
\xlink{Gpg4win for Advanced Users}{durchblicker.html}" 
for additional information on adding signatures to \Email{}s, as well 
as importing and using an existing private key in GnuPG.
}

\textbf{
\OtherBook{} Furthermore, Chapters 13 and 14 of the manual 
"\xlink{Gpg4win for Advanced Users}{durchblicker.html}" give you 
additional information about the ideas behind GnuPG's security
features, as well as detailed insights into the mathematical
concepts used.  These chapters are easy to read, they were
written for \textbf{the every day user}, not mathematicians and
cryptographers. }


\newpage
\appendix
\section{Suggestions regarding the Outlook plugin \emph{GPGol}}

GPGol, a Microsoft Outlook plugin, integrates the operation of GnuPG.
Here are some tips about the operation of this plugin.

While it is fairly easy to operate the plugin in other \Email{} programs,
technical difficulties associated with its integration of OpenPGP
into Outlook make it harder to operate in this context.

The current version of GPGol works only with Outlook 2003 SP2 and will
notify you if you are using an older version.

Please be aware of the following limitations:

\begin{itemize}
\item Do not use Word as your text editor. 

\item Inline-PGP or the conventional PGP is fully supported, however,
PGP/MIME \Email{}s cannot be created. 

On the other hand, decryption and signature verification of simple,
as opposed to more complex, PGP/MIME \Email{}s, is not supported.
\end{itemize}


\subsection{Installation}

The plugin can be installed using the the Gpg4win Installer. The next
time you start Outlook, a \Menu{GnuPG} tab will be added in the
\Menu{Tools$\rightarrow$Options} menu :

\begin{center}
\IncludeImage[width=0.7\textwidth]{sc-gpgol-options}
\end{center}

The first two options of this tab allow you to set default settings
to do with encrypting and signing new messages.  Note that you
can always change these settings when creating new messages.

The option "Save decrypted message automatically" will save attachments
as unencrypted text (after decryption). As a result, the attachment can
be read anytime without having to decrypt it again. 

The option "Automatically sign attachments" ensures that attachments are
automatically signed along with your main text. For this purpose, an
attachment with a signature is created for every attachment included in
the message.

The option "Also encrypt message with the default key" allows you to
enter the Key ID of your own key so that your messages are automatically
(and additionally) encrypted using your own key. This allows you to
later retrieve, decrypt and read messages in the Sent folder.

The option "Also decrypt in preview window" is only feasible with very
fast computers; at the present time it has only limited operational
capacity. 

The option "Show HTML view if possible" enables you to view a message
in HTML format. In most cases, or when an HTML format is not available,
it will be shown in text format. 

There are a number of preset default settings which take effect after
initial installation. However, you should also ensure that you are
not using \textbf{Microsoft Word} to compose messages.  We also do not
recommend using HTML messages.

You can set these option using the \Menu{Mail Format} tab in the
\Menu{Tools$\rightarrow$Options} menu.  Your settings should look like
this:

\begin{center}
\IncludeImage[width=0.7\textwidth]{sc-gpgol-noword}
\end{center}


\subsection{Common Questions}

\begin{description}
\item[Encryption Icons are not visible in the notification area/system
tray.] 

%The german text below was somewhat confusing - I know what they are
%trying to accomplish here and changed the english text slightly.

If the system tray contains a large number of icons, it may not display
all of them. Clicking on the arrow at the right end of the system
tray (located on the bottom right hand side of most computers) will open
up the tray to show all icons, including the Key Manager icon.

Again, make sure that you are not using Microsoft Word to compose
messages (see installation instructions).
 
\item[What is the meaning of the letters included in the GnuPG tab?] 

  GPGol lists the components and their associated file names as part of
the decryption or message signature verification. File names
containing encrypted attachments are prefaced with "E", and signed
attachments are marked with an "S". 

\item[Where can I locate information regarding the latest version
GPGol?]       

  Click on the logo on the lower left side of your GnuPG options menu.

\item[Why is it not possible to cancel an encryption process?] 

  GPGol begins the encryption process as soon as you press the Send
button. An error contained in Outlook prevents you from interrupting or
cancelling this process. To avoid this, we recommend configuring Outlook
so that messages are not sent immediately. This gives you the
opportunity to cancel the message before it is sent. As a security
measure, GPGol will try to erase the contents of the message if the
encryption process is interrupted, but this is not always successful.

\item[Why does the confirmation window appear when accessing certain
\Email{}s?]

If \linebreak{} GPGol is not installed as a trusted plugin, Outlook
assumes that it is trying to access internal information without
authorization.  GPGol attempts to prevent this from happening, but
sometimes it is necessary in order to properly display encrypted or
signed \Email{}s.

GPGol is still "in development". One outstanding issue is its
registration as a trusted plugin. Future versions of GPGol should be able
to solve the issue of needing to confirm some or all \Email{}s.

\item[How come GPGol cannot create PGP/MIME messages?]

At the present, there is no way of telling Outlook that a PGP/MIME
message is to be created. Outlook decides on its own which
"Content-Type" is to be used, as the plugin is not able to preset a
certain type.  You may want to contact Microsoft directly to talk about
this setting and/or missing documentation regarding this setup.

\item[Why are signature verifications not performed automatically?]

We are currently working on an automatic signature verification option
to be used when opening an \Email{}.  However, given Outlook's complex
technical environment, this is not an easy issue to solve. 


\end{description}

\clearpage
\section{Transferring from other GnuPG programs}
\label{ch:migration}

This section explains how to migrate from out GnuPG-based programs to Gpg4win.
The installation program recognizes some of these programs
and notifies you if that is the case.

As a general rule it is recommended that existing GnuPG-based programs
be removed before installing Gpg4win. 

\textbf{Always} save a backup of the existing keys. 
The best way to do this is to use
the options in the existing system before installation Gpg4win.
Select the option to save any private (secret) keys, 
as well as all existing public keys.  Save these in one or two files.

As soon as Gpg4win is installed, check if your existing keys are listed
in the new program (using either GPA or WinPT). If yes, Gpg4win 
was able to detect the previous key storage and you
need take no further action.

If the program does not list the existing keys, you must import them from
your backup file(s).  For more information on this topic, refer to Chapter
12 in the manual "\xlink{Gpg4win for Avanced
Users}{durchblicker.html}".

If your older system also includes GPA, you can use its back-up option,
which should be very similar to the operation of GPA in Gpg4win.

If you are not able to locate your existing keys, use the 'Search'
function in Windows to find the files named \Filename{secring.gpg} and
\Filename{pubring.gpg} an import them per GPA.\footnote{This is not the
official procedure, but it works with all current versions of GnuPG.}




\clearpage
\section{History}

\begin{itemize}
\item  "GnuPP for Beginners", First Edition March 2002,\\
  Authors: Manfred J. Heinze, TextLab text+media\\
  Consulting: Lutz Zolondz, G-N-U GmbH\\
  Illustrations: Karl Bihlmeier, Bihlmeier \& Kramer GbR\\
  Layout: Isabel Kramer, Bihlmeier \& Kramer GbR\\
  Documentation: Dr. Francis Wray, e-mediate Ltd.\\
  Editor: Ute Bahn, TextLab text+media\\
  Published by the German Federal Ministry for Industry \& Technology.\\
  Available at
  \verb-http://www.gnupp.de/pdf/einsteiger.pdf-.
%    Der Abschnitt "History" ist im Originaldokument nicht vorhanden
%    und wurde von Werner Koch beigef�gt.
\item Revised unpublished version by TextLab text+media.      
\item "Gpg4win for Beginners", December 2005\\
      Authors:  Werner Koch, g10 Code GmbH\\
      Published as part of the Gpg4win project.
\item "Gpg4win for Novices", November 2006\\
      Translated by: Brigitte Hamilton\\
      Published as part of the Gpg4win project.

\end{itemize}

\clearpage
\section{GNU Free Documentation License}

 \begin{center}

       Version 1.2, November 2002


 Copyright \copyright 2000,2001,2002  Free Software Foundation, Inc.
 
 \bigskip
 
     51 Franklin St, Fifth Floor, Boston, MA  02110-1301  USA
  
 \bigskip
 
 Everyone is permitted to copy and distribute verbatim copies
 of this license document, but changing it is not allowed.
\end{center}


\begin{center}
{\bf\large Preamble}
\end{center}

The purpose of this License is to make a manual, textbook, or other
functional and useful document ``free'' in the sense of freedom: to
assure everyone the effective freedom to copy and redistribute it,
with or without modifying it, either commercially or noncommercially.
Secondarily, this License preserves for the author and publisher a way
to get credit for their work, while not being considered responsible
for modifications made by others.

This License is a kind of ``copyleft'', which means that derivative
works of the document must themselves be free in the same sense.  It
complements the GNU General Public License, which is a copyleft
license designed for free software.

We have designed this License in order to use it for manuals for free
software, because free software needs free documentation: a free
program should come with manuals providing the same freedoms that the
software does.  But this License is not limited to software manuals;
it can be used for any textual work, regardless of subject matter or
whether it is published as a printed book.  We recommend this License
principally for works whose purpose is instruction or reference.


\begin{center}
{\Large\bf 1. APPLICABILITY AND DEFINITIONS}
%\addcontentsline{toc}{section}{1. APPLICABILITY AND DEFINITIONS}
\end{center}

This License applies to any manual or other work, in any medium, that
contains a notice placed by the copyright holder saying it can be
distributed under the terms of this License.  Such a notice grants a
world-wide, royalty-free license, unlimited in duration, to use that
work under the conditions stated herein.  The \textbf{``Document''}, below,
refers to any such manual or work.  Any member of the public is a
licensee, and is addressed as \textbf{``you''}.  You accept the license if you
copy, modify or distribute the work in a way requiring permission
under copyright law.

A \textbf{``Modified Version''} of the Document means any work containing the
Document or a portion of it, either copied verbatim, or with
modifications and/or translated into another language.

A \textbf{``Secondary Section''} is a named appendix or a front-matter section of
the Document that deals exclusively with the relationship of the
publishers or authors of the Document to the Document's overall subject
(or to related matters) and contains nothing that could fall directly
within that overall subject.  (Thus, if the Document is in part a
textbook of mathematics, a Secondary Section may not explain any
mathematics.)  The relationship could be a matter of historical
connection with the subject or with related matters, or of legal,
commercial, philosophical, ethical or political position regarding
them.

The \textbf{``Invariant Sections''} are certain Secondary Sections whose titles
are designated, as being those of Invariant Sections, in the notice
that says that the Document is released under this License.  If a
section does not fit the above definition of Secondary then it is not
allowed to be designated as Invariant.  The Document may contain zero
Invariant Sections.  If the Document does not identify any Invariant
Sections then there are none.

The \textbf{``Cover Texts''} are certain short passages of text that are listed,
as Front-Cover Texts or Back-Cover Texts, in the notice that says that
the Document is released under this License.  A Front-Cover Text may
be at most 5 words, and a Back-Cover Text may be at most 25 words.

A \textbf{``Transparent''} copy of the Document means a machine-readable copy,
represented in a format whose specification is available to the
general public, that is suitable for revising the document
straightforwardly with generic text editors or (for images composed of
pixels) generic paint programs or (for drawings) some widely available
drawing editor, and that is suitable for input to text formatters or
for automatic translation to a variety of formats suitable for input
to text formatters.  A copy made in an otherwise Transparent file
format whose markup, or absence of markup, has been arranged to thwart
or discourage subsequent modification by readers is not Transparent.
An image format is not Transparent if used for any substantial amount
of text.  A copy that is not ``Transparent'' is called \textbf{``Opaque''}.

Examples of suitable formats for Transparent copies include plain
ASCII without markup, Texinfo input format, LaTeX input format, SGML
or XML using a publicly available DTD, and standard-conforming simple
HTML, PostScript or PDF designed for human modification.  Examples of
transparent image formats include PNG, XCF and JPG.  Opaque formats
include proprietary formats that can be read and edited only by
proprietary word processors, SGML or XML for which the DTD and/or
processing tools are not generally available, and the
machine-generated HTML, PostScript or PDF produced by some word
processors for output purposes only.

The \textbf{``Title Page''} means, for a printed book, the title page itself,
plus such following pages as are needed to hold, legibly, the material
this License requires to appear in the title page.  For works in
formats which do not have any title page as such, ``Title Page'' means
the text near the most prominent appearance of the work's title,
preceding the beginning of the body of the text.

A section \textbf{``Entitled XYZ''} means a named subunit of the Document whose
title either is precisely XYZ or contains XYZ in parentheses following
text that translates XYZ in another language.  (Here XYZ stands for a
specific section name mentioned below, such as \textbf{``Acknowledgements''},
\textbf{``Dedications''}, \textbf{``Endorsements''}, or \textbf{``History''}.)  
To \textbf{``Preserve the Title''}
of such a section when you modify the Document means that it remains a
section ``Entitled XYZ'' according to this definition.

The Document may include Warranty Disclaimers next to the notice which
states that this License applies to the Document.  These Warranty
Disclaimers are considered to be included by reference in this
License, but only as regards disclaiming warranties: any other
implication that these Warranty Disclaimers may have is void and has
no effect on the meaning of this License.


\begin{center}
{\Large\bf 2. VERBATIM COPYING}
%\addcontentsline{toc}{section}{2. VERBATIM COPYING}
\end{center}

You may copy and distribute the Document in any medium, either
commercially or noncommercially, provided that this License, the
copyright notices, and the license notice saying this License applies
to the Document are reproduced in all copies, and that you add no other
conditions whatsoever to those of this License.  You may not use
technical measures to obstruct or control the reading or further
copying of the copies you make or distribute.  However, you may accept
compensation in exchange for copies.  If you distribute a large enough
number of copies you must also follow the conditions in section 3.

You may also lend copies, under the same conditions stated above, and
you may publicly display copies.


\begin{center}
{\Large\bf 3. COPYING IN QUANTITY}
%\addcontentsline{toc}{section}{3. COPYING IN QUANTITY}
\end{center}


If you publish printed copies (or copies in media that commonly have
printed covers) of the Document, numbering more than 100, and the
Document's license notice requires Cover Texts, you must enclose the
copies in covers that carry, clearly and legibly, all these Cover
Texts: Front-Cover Texts on the front cover, and Back-Cover Texts on
the back cover.  Both covers must also clearly and legibly identify
you as the publisher of these copies.  The front cover must present
the full title with all words of the title equally prominent and
visible.  You may add other material on the covers in addition.
Copying with changes limited to the covers, as long as they preserve
the title of the Document and satisfy these conditions, can be treated
as verbatim copying in other respects.

If the required texts for either cover are too voluminous to fit
legibly, you should put the first ones listed (as many as fit
reasonably) on the actual cover, and continue the rest onto adjacent
pages.

If you publish or distribute Opaque copies of the Document numbering
more than 100, you must either include a machine-readable Transparent
copy along with each Opaque copy, or state in or with each Opaque copy
a computer-network location from which the general network-using
public has access to download using public-standard network protocols
a complete Transparent copy of the Document, free of added material.
If you use the latter option, you must take reasonably prudent steps,
when you begin distribution of Opaque copies in quantity, to ensure
that this Transparent copy will remain thus accessible at the stated
location until at least one year after the last time you distribute an
Opaque copy (directly or through your agents or retailers) of that
edition to the public.

It is requested, but not required, that you contact the authors of the
Document well before redistributing any large number of copies, to give
them a chance to provide you with an updated version of the Document.


\begin{center}
{\Large\bf 4. MODIFICATIONS}
%\addcontentsline{toc}{section}{4. MODIFICATIONS}
\end{center}

You may copy and distribute a Modified Version of the Document under
the conditions of sections 2 and 3 above, provided that you release
the Modified Version under precisely this License, with the Modified
Version filling the role of the Document, thus licensing distribution
and modification of the Modified Version to whoever possesses a copy
of it.  In addition, you must do these things in the Modified Version:

\begin{itemize}
\item[A.] 
   Use in the Title Page (and on the covers, if any) a title distinct
   from that of the Document, and from those of previous versions
   (which should, if there were any, be listed in the History section
   of the Document).  You may use the same title as a previous version
   if the original publisher of that version gives permission.
   
\item[B.]
   List on the Title Page, as authors, one or more persons or entities
   responsible for authorship of the modifications in the Modified
   Version, together with at least five of the principal authors of the
   Document (all of its principal authors, if it has fewer than five),
   unless they release you from this requirement.
   
\item[C.]
   State on the Title page the name of the publisher of the
   Modified Version, as the publisher.
   
\item[D.]
   Preserve all the copyright notices of the Document.
   
\item[E.]
   Add an appropriate copyright notice for your modifications
   adjacent to the other copyright notices.
   
\item[F.]
   Include, immediately after the copyright notices, a license notice
   giving the public permission to use the Modified Version under the
   terms of this License, in the form shown in the Addendum below.
   
\item[G.]
   Preserve in that license notice the full lists of Invariant Sections
   and required Cover Texts given in the Document's license notice.
   
\item[H.]
   Include an unaltered copy of this License.
   
\item[I.]
   Preserve the section Entitled ``History'', Preserve its Title, and add
   to it an item stating at least the title, year, new authors, and
   publisher of the Modified Version as given on the Title Page.  If
   there is no section Entitled ``History'' in the Document, create one
   stating the title, year, authors, and publisher of the Document as
   given on its Title Page, then add an item describing the Modified
   Version as stated in the previous sentence.
   
\item[J.]
   Preserve the network location, if any, given in the Document for
   public access to a Transparent copy of the Document, and likewise
   the network locations given in the Document for previous versions
   it was based on.  These may be placed in the ``History'' section.
   You may omit a network location for a work that was published at
   least four years before the Document itself, or if the original
   publisher of the version it refers to gives permission.
   
\item[K.]
   For any section Entitled ``Acknowledgements'' or ``Dedications'',
   Preserve the Title of the section, and preserve in the section all
   the substance and tone of each of the contributor acknowledgements
   and/or dedications given therein.
   
\item[L.]
   Preserve all the Invariant Sections of the Document,
   unaltered in their text and in their titles.  Section numbers
   or the equivalent are not considered part of the section titles.
   
\item[M.]
   Delete any section Entitled ``Endorsements''.  Such a section
   may not be included in the Modified Version.
   
\item[N.]
   Do not retitle any existing section to be Entitled ``Endorsements''
   or to conflict in title with any Invariant Section.
   
\item[O.]
   Preserve any Warranty Disclaimers.
\end{itemize}

If the Modified Version includes new front-matter sections or
appendices that qualify as Secondary Sections and contain no material
copied from the Document, you may at your option designate some or all
of these sections as invariant.  To do this, add their titles to the
list of Invariant Sections in the Modified Version's license notice.
These titles must be distinct from any other section titles.

You may add a section Entitled ``Endorsements'', provided it contains
nothing but endorsements of your Modified Version by various
parties--for example, statements of peer review or that the text has
been approved by an organization as the authoritative definition of a
standard.

You may add a passage of up to five words as a Front-Cover Text, and a
passage of up to 25 words as a Back-Cover Text, to the end of the list
of Cover Texts in the Modified Version.  Only one passage of
Front-Cover Text and one of Back-Cover Text may be added by (or
through arrangements made by) any one entity.  If the Document already
includes a cover text for the same cover, previously added by you or
by arrangement made by the same entity you are acting on behalf of,
you may not add another; but you may replace the old one, on explicit
permission from the previous publisher that added the old one.

The author(s) and publisher(s) of the Document do not by this License
give permission to use their names for publicity for or to assert or
imply endorsement of any Modified Version.


\begin{center}
{\Large\bf 5. COMBINING DOCUMENTS}
%\addcontentsline{toc}{section}{5. COMBINING DOCUMENTS}
\end{center}


You may combine the Document with other documents released under this
License, under the terms defined in section 4 above for modified
versions, provided that you include in the combination all of the
Invariant Sections of all of the original documents, unmodified, and
list them all as Invariant Sections of your combined work in its
license notice, and that you preserve all their Warranty Disclaimers.

The combined work need only contain one copy of this License, and
multiple identical Invariant Sections may be replaced with a single
copy.  If there are multiple Invariant Sections with the same name but
different contents, make the title of each such section unique by
adding at the end of it, in parentheses, the name of the original
author or publisher of that section if known, or else a unique number.
Make the same adjustment to the section titles in the list of
Invariant Sections in the license notice of the combined work.

In the combination, you must combine any sections Entitled ``History''
in the various original documents, forming one section Entitled
``History''; likewise combine any sections Entitled ``Acknowledgements'',
and any sections Entitled ``Dedications''.  You must delete all sections
Entitled ``Endorsements''.

\begin{center}
{\Large\bf 6. COLLECTIONS OF DOCUMENTS}
%\addcontentsline{toc}{section}{6. COLLECTIONS OF DOCUMENTS}
\end{center}

You may make a collection consisting of the Document and other documents
released under this License, and replace the individual copies of this
License in the various documents with a single copy that is included in
the collection, provided that you follow the rules of this License for
verbatim copying of each of the documents in all other respects.

You may extract a single document from such a collection, and distribute
it individually under this License, provided you insert a copy of this
License into the extracted document, and follow this License in all
other respects regarding verbatim copying of that document.


\begin{center}
{\Large\bf 7. AGGREGATION WITH INDEPENDENT WORKS}
%\addcontentsline{toc}{section}{7. AGGREGATION WITH INDEPENDENT WORKS}
\end{center}


A compilation of the Document or its derivatives with other separate
and independent documents or works, in or on a volume of a storage or
distribution medium, is called an ``aggregate'' if the copyright
resulting from the compilation is not used to limit the legal rights
of the compilation's users beyond what the individual works permit.
When the Document is included in an aggregate, this License does not
apply to the other works in the aggregate which are not themselves
derivative works of the Document.

If the Cover Text requirement of section 3 is applicable to these
copies of the Document, then if the Document is less than one half of
the entire aggregate, the Document's Cover Texts may be placed on
covers that bracket the Document within the aggregate, or the
electronic equivalent of covers if the Document is in electronic form.
Otherwise they must appear on printed covers that bracket the whole
aggregate.


\begin{center}
{\Large\bf 8. TRANSLATION}
%\addcontentsline{toc}{section}{8. TRANSLATION}
\end{center}


Translation is considered a kind of modification, so you may
distribute translations of the Document under the terms of section 4.
Replacing Invariant Sections with translations requires special
permission from their copyright holders, but you may include
translations of some or all Invariant Sections in addition to the
original versions of these Invariant Sections.  You may include a
translation of this License, and all the license notices in the
Document, and any Warranty Disclaimers, provided that you also include
the original English version of this License and the original versions
of those notices and disclaimers.  In case of a disagreement between
the translation and the original version of this License or a notice
or disclaimer, the original version will prevail.

If a section in the Document is Entitled ``Acknowledgements'',
``Dedications'', or ``History'', the requirement (section 4) to Preserve
its Title (section 1) will typically require changing the actual
title.


\begin{center}
{\Large\bf 9. TERMINATION}
%\addcontentsline{toc}{section}{9. TERMINATION}
\end{center}


You may not copy, modify, sublicense, or distribute the Document except
as expressly provided for under this License.  Any other attempt to
copy, modify, sublicense or distribute the Document is void, and will
automatically terminate your rights under this License.  However,
parties who have received copies, or rights, from you under this
License will not have their licenses terminated so long as such
parties remain in full compliance.


\begin{center}
{\Large\bf 10. FUTURE REVISIONS OF THIS LICENSE}
%%\addcontentsline{toc}{section}{10. FUTURE REVISIONS OF THIS LICENSE}
\end{center}


The Free Software Foundation may publish new, revised versions
of the GNU Free Documentation License from time to time.  Such new
versions will be similar in spirit to the present version, but may
differ in detail to address new problems or concerns.  See
http://www.gnu.org/copyleft/.

Each version of the License is given a distinguishing version number.
If the Document specifies that a particular numbered version of this
License ``or any later version'' applies to it, you have the option of
following the terms and conditions either of that specified version or
of any later version that has been published (not as a draft) by the
Free Software Foundation.  If the Document does not specify a version
number of this License, you may choose any version ever published (not
as a draft) by the Free Software Foundation.


\begin{center}
{\Large\bf ADDENDUM: How to use this License for your documents}
%\addcontentsline{toc}{section}{ADDENDUM: How to use this License for your documents}
\end{center}

To use this License in a document you have written, include a copy of
the License in the document and put the following copyright and
license notices just after the title page:

\bigskip
\begin{quote}
    Copyright \copyright  YEAR  YOUR NAME.
    Permission is granted to copy, distribute and/or modify this document
    under the terms of the GNU Free Documentation License, Version 1.2
    or any later version published by the Free Software Foundation;
    with no Invariant Sections, no Front-Cover Texts, and no Back-Cover Texts.
    A copy of the license is included in the section entitled ``GNU
    Free Documentation License''.
\end{quote}
\bigskip
    
If you have Invariant Sections, Front-Cover Texts and Back-Cover Texts,
replace the ``with...Texts.'' line with this:

\bigskip
\begin{quote}
    with the Invariant Sections being LIST THEIR TITLES, with the
    Front-Cover Texts being LIST, and with the Back-Cover Texts being LIST.
\end{quote}
\bigskip
    
If you have Invariant Sections without Cover Texts, or some other
combination of the three, merge those two alternatives to suit the
situation.

If your document contains nontrivial examples of program code, we
recommend releasing these examples in parallel under your choice of
free software license, such as the GNU General Public License,
to permit their use in free software.



\end{document}
